\section{Flow and Vorticity}
\label{sec:flow}

\subsection{Log-Flow}

The antisymmetric component of $C_{ij}$ defines the \textbf{log-flow}:

\begin{equation}
  J_{ij} = C_{ij} - C_{ji} = 2A_{ij}.
  \label{eq:flow}
\end{equation}

$J_{ij}$ represents the net directional transfer of causal influence.
It satisfies $J_{ij} = -J_{ji}$ by construction.
Motion arises from asymmetry: where $J_{ij} = 0$, no net transfer occurs;
where $J_{ij} \neq 0$, directed propagation occurs.

\subsection{The Continuity Equation}

We use $\mathbf{J}(x)$ as a coarse-grained continuum representation of the
discrete log-flow $J_{ij}$. This is an informal bridge: an explicit averaging,
scaling, orientation, and discrete-to-vector-field map is not supplied here.
The evolution of continuum log density $\rho$ is written as:

\begin{keybox}
\begin{equation}
  \partial_t\rho + \nabla \cdot \mathbf{J} = \sigma,
  \label{eq:continuity}
\end{equation}
\end{keybox}

Equation~\eqref{eq:continuity} is a balance / continuity relation within this
continuum field representation. Because the source term may be nonzero, the
display is not an unconditional conservation law and does not, by its balance
form alone, establish a physical-time interpretation.

The continuum field $\rho$ is associated with the node quantity $\rho_i$, but
no interpolation, volume normalization, unit map, or complete function
signature relating them is supplied here. Both continuum bridges remain open.

Here $\sigma \geq 0$ is the \textbf{log-generation rate}.
$\sigma$ is not an external source term.
It is the rate at which the system fails to remain closed:
wherever structure is incomplete, new logs are generated.

\begin{remark}[Source-local symbol]
Here $\sigma$ denotes the Vol.~1 continuum log-generation rate. It is not
identified by notation alone with the microscopic $\dot{\rho}_i$ expression
or with IFGT's separately constructed generation term $\sigma$; those
translations remain unspecified.
\end{remark}

This is a balance relation rather than an unconditional conservation law.
Within the stated continuum model, $\sigma = 0$ marks the source-free regime
associated with the classical fixed-point phase (\cref{sec:phase}).
This formal regime distinction does not by itself establish whether physical
conservation laws are fundamental or emergent.

\begin{keybox}
\begin{equation}
  \sigma > 0
  \qquad\text{(positive source regime).}
  \label{eq:sigma_time}
\end{equation}
\end{keybox}

Within this source-local continuum model, $\sigma>0$ denotes positive log
generation and is proposed as an indicator of accumulation / temporal
ordering.  Neither direction of an equivalence with physical-time generation
is established by the balance equation.

\subsection{Helmholtz Decomposition}

The flow field admits a Helmholtz decomposition:

\[
  \mathbf{J} = \mathbf{J}_{\text{rot}} + \mathbf{J}_{\text{pot}},
\]

where $\nabla \cdot \mathbf{J}_{\text{rot}} = 0$
and $\nabla \times \mathbf{J}_{\text{pot}} = 0$.
The solenoidal component $\mathbf{J}_{\text{rot}}$ carries vorticity and is
used in the proposed enclosure mechanism.  Persistence and stability require
additional dynamical analysis.
The irrotational component $\mathbf{J}_{\text{pot}}$ carries the divergence
associated with the log-generation rate $\sigma$.

\subsection{Vorticity and Closure Density}

Define the \textbf{vorticity} of the log-flow:

\begin{equation}
  \omega = \nabla \times \mathbf{J}.
  \label{eq:vorticity}
\end{equation}

This is a definition, and hence a derived construction, in a continuum
vector-field setting in which the curl operator is meaningful. It defines the
vorticity of the represented log-flow; nonzero $\omega$ alone is not a
stability or persistence theorem.

Vorticity measures local circulation.
In the subsequent enclosure interpretation of the model, $\omega=0$ denotes
an irrotational represented flow.  When $\omega\neq0$, circulation is present
and flow lines may close; this is a candidate structural route to enclosure,
not a theorem of persistence.

The local intensity of vortical structure defines the \textbf{closure density}:

\begin{equation}
  L \sim |\omega|.
  \label{eq:closure_density}
\end{equation}

This is an effective proportional identification of a scalar closure-intensity
measure in the present representation. The magnitude convention and
proportionality normalization are not fixed here, and the relation is not a
derived universal identity.

\begin{remark}[Carrier]
The $L$ in this section is a scalar local closure-density / vortical-intensity
quantity. The edge quantity $L_{ab}$ and coarse tensor $L_{\mu\nu}$ introduced
in Vol.~3 are distinct carrier-specific objects; shared lettering does not
define a map among them.
\end{remark}

$L$ is used in Vols.~2--3 as part of candidate structural interpretations of
particle-like and gravitational behavior; those physical identifications are
separate from the definition above.

\subsection{Role-Separated Construction Summary}

The following historical compact sequence summarizes the route from asymmetry
to the scalar closure measure:

\[
  C_{ij}
  \;\xrightarrow{J = C - C^T}\; J
  \;\xrightarrow{\omega = \nabla\times J}\; \omega
  \;\xrightarrow{L \sim |\omega|}\; L.
\]

Its first edge is a discrete definition, followed by an OPEN
discrete-to-continuum representation bridge before the curl definition; the
last edge is an effective proportional relation.  The arrows therefore do not
share one logical status, and the sequence alone does not establish persistent
physical enclosure.
